% =====================================================================
% Revised Methods material for:
%   "Machine learning in battery management systems: a two-tier scoping
%    review of 799 eligible records with in-depth classification and
%    validation audit of 103 studies" (TTE-Reg-2026-10-3720)
%
% Purpose: drop-in replacement / addition for the Methods section in the
% revised manuscript. These paragraphs implement the technical-rigor
% upgrades requested by the Associate Editors:
%   (AE-1) deeper technical treatment; transportation-specific emphasis
%   (AE-2) rigorous independent full-text verification of the validation
%          audit (response to kappa = 0.48), justification of the 103-
%          study selection, treatment of selection bias, EV-specific BMS
%          applications, vehicle-level validation, onboard implementation,
%          transportation safety requirements.
%
% Usage notes:
%  * No numbers are invented here. Results-dependent statements are phrased
%    so they can be completed once the re-audit is finished (e.g. after
%    "the re-audit achieved", insert the computed value from
%    kappa_analysis.py).
%  * % INSERT-CITATION comments mark places where the authors may re-attach
%    references from the original manuscript bibliography.
%  * The reliability protocol assumes the companion instrument package:
%    Validation_Audit_Codebook_v2.docx, Coding_Workbook_v2.xlsx,
%    kappa_analysis.py.
% =====================================================================

% ---------------------------------------------------------------------
% 1. RELIABILITY OF THE VALIDATION AUDIT  (replaces the single paragraph
%    that previously reported kappa = 0.48)
% ---------------------------------------------------------------------
\subsection{Reliability of the validation audit}

The initial pilot audit of validation practices, in which two authors
independently coded a sample of studies, yielded a Cohen's
$\kappa$ of 0.48 for the validation-protocol fields. We interpreted this
as evidence that the version~1 instrument, rather than the underlying
literature, was under-specified: post-hoc inspection of the
disagreements showed that the majority arose from (i) adjacent-category
confusion on ordinal ladders (notably the boundary between
same-cell/different-condition and different-cell testing), (ii)
undefined treatment of multi-experiment studies, and (iii) silent
defaults for unreported items. We therefore rebuilt the instrument
before drawing any conclusion that depends on it.

The validation audit was repeated with a version~2.0 codebook
(supplementary material) that specifies, for each of fourteen variables
(V1--V14): the construct, an exhaustive category set, decision rules
for multi-experiment and multi-target studies, and worked examples for
every known edge case. Three design changes target the failure modes of
the pilot: ordinal ladders were re-anchored with demonstrable-evidence
tests (an item is coded at a level only if the paper reports results at
that level, not because the authors claim it); ``not reported'' became
an explicit, admissible code so that coders never infer silently; and
any code at the two highest levels of an ordinal scale now requires a
verbatim quotation with a section or page reference in the evidence
column.

Two coders (C1, C2), both co-authors, independently re-coded all 103
studies from the full text and supplementary material. Coder~1 had
participated in the original audit; coder~2 had not. Both were blind to
the original codes and to each other's sheets. Coding proceeded in two
stages. First, a calibration round on twelve studies stratified across
publication years and transportation-relevance strata; the
pre-registered gate for proceeding was per-variable $\kappa \geq 0.70$
on the calibration set. Discrepancies in the calibration round were
resolved by revising the codebook (three iterations), not by
negotiating codes. Second, the full double-coding of all 103 studies;
every remaining disagreement was adjudicated by a third senior author
(A) within 48 hours, and every resolution was logged with a one-line
rationale and the codebook clause on which it rests. The adjudication
log is published as a supplementary file.

Agreement is reported per variable as raw percent agreement, Cohen's
$\kappa$, and PABAK (prevalence-adjusted bias-adjusted kappa), because
several categories are rare and prevalence-dependent; for the six
ordinal variables we additionally report linearly and quadratically
weighted $\kappa$. Uncertainty is quantified with bootstrap 95\%
confidence intervals (10{,}000 resamples, percentile method). We
pre-declared a category-collapse rule: if a single pair of adjacent
categories accounts for more than half of the disagreements on a
variable, we report both the granular and the collapsed $\kappa$ and
use the collapsed variable in sensitivity analyses. The initial pilot
value ($\kappa = 0.48$) is retained in the manuscript alongside the
re-audited values, so that the improvement in the measurement
instrument, and its effect on the findings, is transparent.
% INSERT-CITATION: Cohen 1960; Fleiss 1971; Landis & Koch 1977; Byrt et al. 1993

The re-audit achieved [COMPLETE AFTER RE-CODING: per-variable kappa,
PABAK, weighted kappa, and bootstrap CIs from kappa\_analysis.py].
Variables that remained below $\kappa = 0.60$ after codebook revision
are dichotomised for all headline analyses, and the granular
distributions are reported in the appendix; no variable is silently
dropped. All conclusions that rest on validation-audit variables are
recomputed from the adjudicated dataset, and where the re-audit changes
a finding reported in the original submission, the change is stated
explicitly.

% ---------------------------------------------------------------------
% 2. SELECTION OF THE 103 STUDIES  (new subsection)
% ---------------------------------------------------------------------
\subsection{Selection of studies for the in-depth audit}

The 103 studies audited in depth (tier~2) were drawn from the 799
eligible records by explicit, pre-declared criteria rather than by
convenience. A record entered tier~2 if it (a) reports an original
machine-learning method or an empirical evaluation of one for a battery
management task (thus excluding reviews, surveys, and purely
position papers), and (b) reports at least one quantitative result on
battery data, so that validation practices are observable in principle.
Records meeting (a) but not (b) --- for example, architecture proposals
without experiments --- were retained in tier~1 but excluded from the
audit, and their count and characteristics are reported so that the
denominator is reconstructable at every step.

Because the associate editors correctly noted that the selection of the
103 studies requires stronger justification and a clearer treatment of
selection bias, the revision adds three safeguards. First, a complete
PRISMA-ScR flow diagram with counts and reasons at every exclusion
step, replacing the summary numbers previously given in the text.
Second, a characterization of the excluded-but-eligible records (venue,
year, task, and validation keywords) demonstrating whether the audited
set differs systematically from the records it was drawn from.
Third, a sensitivity analysis in which every headline finding is
recomputed on the subset of tier-2 studies that passed the selection
criteria unambiguously (no borderline judgement involved), so that
readers can verify that the findings are not an artefact of the
inclusion rule. Records whose inclusion involved a borderline judgement
are listed individually in the appendix with the reason for inclusion.
% INSERT-CITATION: Tricco et al. 2018 (PRISMA-ScR)

% ---------------------------------------------------------------------
% 3. RECALL ASSURANCE  (new subsection; answers the recall-audit
%    criticism)
% ---------------------------------------------------------------------
\subsection{Search recall assurance}

Our own recall audit indicated that the search strategy did not
retrieve a portion of the eligible literature. The revision
strengthens retrieval assurance in four ways. (i)~Backward and forward
citation chaining on all 103 included studies (forward chaining via
Google Scholar and OpenAlex citation indexes), with newly identified
records screened against the same eligibility criteria. (ii)~A
hand-search of the tables of contents of five venues in which eligible
records concentrate --- Journal of Power Sources, eTransportation,
Journal of Energy Storage, and the IEEE transportation-electrification
venues --- over the search window. (iii)~A 30-record benchmark set of
studies that the co-authors independently judged must be found by any
adequate search (``known-must-find'' set), assembled before re-running
the searches, against which recall is measured explicitly. (iv)~A
re-executed search with the full, dated query strings, databases, and
filters documented in a provenance appendix, so the search is
reproducible. Recall against the benchmark set, and the number of
additional eligible records recovered by chaining and hand-search, are
reported as a range with the audit trail, and the tier-2 selection
described above is re-run if the recovered records change the eligible
pool materially.

% ---------------------------------------------------------------------
% 4. TRANSPORTATION-SPECIFIC SYNTHESIS  (new subsection; answers AE-1
%    and the EV emphasis of AE-2)
% ---------------------------------------------------------------------
\subsection{Transportation-specific synthesis}

Because the broader battery literature otherwise dominates the
analysis, the revision adds a dedicated synthesis of the
transportation-facing subset, coded with three new variables. Variable
V11 (deployment-technology-readiness ladder) distinguishes offline
laboratory validation (level~0), dynamic-profile bench testing with
drive-cycle current profiles on a cycler (level~1), hardware-in-the-loop
or emulated BMS hardware (level~2), retrospective or live data from
real vehicles and fleets (level~3), and demonstrated onboard deployment
(level~4). Variable V12 records whether onboard feasibility is
addressed at all: not discussed, mentioned qualitatively, quantified
against a latency / memory / compute budget, or measured on target
hardware. Variable V14 separates studies whose validation context is
genuinely vehicle-oriented (EV/HEV, or other transport modes) from
generic battery studies regardless of title framing; coding follows the
validation context, not the claims in the title or introduction.

These variables feed a transportation-specific analysis with four
components. (i)~The distribution of V11/V12 levels across the audited
studies, quantifying how much of the ML-BMS literature has been tested
beyond static laboratory data --- the central transportation-safety
question. (ii)~A vehicle-level validation synthesis restricted to
level~2--4 studies, covering drive-cycle realism (e.g., WLTP, UDDS,
HWFET-type profiles), hardware-in-the-loop setups, and vehicle or
fleet telemetry, with the validation gaps of each study catalogued.
(iii)~An onboard-implementability discussion for the subset of methods
with quantified compute budgets, relating reported footprints to
typical BMS electronic control unit constraints. (iv)~A
transportation-safety framing that connects the observed validation
shortfall to safety-relevant failure modes of learned estimators
(sensor bias and drift, thermal-runaway early warning, state-estimation
error propagation into power limits) and to the assurance expectations
that safety-engineered components face in automotive practice, e.g.,
the argumentation style of ISO~26262; where ML components cannot yet
satisfy such argumentation, we state this as a research gap rather than
claiming readiness.

% ---------------------------------------------------------------------
% 5. OPTIONAL: one-sentence Methods index of the instrument package
% ---------------------------------------------------------------------
% The validation-audit instrument (codebook v2.0, coding workbook, and
% the agreement-analysis script) is released as supplementary material
% and in the project repository, so that the audit is fully replicable.
